% Part: proof-theory
% Chapter: normalization
% Section: introduction

\documentclass[../../../include/open-logic-section]{subfiles}

\begin{document}

\olfileid{pt}{nor}{int}

\olsection{పరిచయం}

షరతు వాక్యం~$!A \lif !B$ను మీరు \tetoken{నిరూపించి}{!!{prove}} ఉంటే,~\Elim{\lif} ద్వారా దాన్ని~$!B$ యొక్క \tetoken{ఒక నిరూపణలో}{!!a{proof}} ఉపయోగించవచ్చు. దీనికి~$!A$ యొక్క \tetoken{ఒక నిరూపణ}{!!a{proof}}~$\delta_1$ కూడా కావాలి. $!A \lif !B$కు మీ \tetoken{నిరూపణ}{!!{proof}} సాధారణంగా~\Intro{\lif}తో ముగుస్తుంది; అది తెరిచి ఉన్న ఉపపత్తి~$!A$ నుంచి~$!B$కు గల \tetoken{ఒక నిరూపణను}{!!a{proof}}~$\delta_2$,~$!A \lif !B$ నిరూపణగా మారుస్తుంది. అలా అయితే మీ చివరి \tetoken{నిరూపణ}{!!{proof}} రూపం ఇది:
\[
\AxiomC{$\Discharge{!A}{x}$}
\RightLabel{$\delta_2$}
\DeduceC{$!B$}
\DischargeRule{\Intro{\lif}}{x}
\UnaryInfC{$!A \lif !B$}
\AxiomC{}
\RightLabel{$\delta_1$}
\DeduceC{$!A$}
\RightLabel{\Elim{\lif}}
\BinaryInfC{$!B$}
\DisplayProof
\]
ఇప్పటికే ఉన్న $!A \lif !B$ యొక్క \tetoken{ఒక నిరూపణను}{!!a{proof}} మీరు మళ్లీ ఉపయోగించడంలో పద్ధతి సమర్థవంతమే; కానీ ఫలిత \tetoken{నిరూపణ}{!!{proof}} అలా కాదు. $!A \lif !B$ను ప్రవేశపెట్టి వెంటనే తొలగించడం పక్కదారిలా కనిపిస్తుంది. $!B$ను \tetoken{నిరూపించడానికి}{!!{prove}} మరింత ప్రత్యక్ష మార్గం ఉండాలి.
\sourcecorrection{OLTEPTNORINT-001}{మూల గద్యంలో $\delta_1$ను ఒకసారి $!A$ నిరూపణగా, మళ్లీ $!A \lif !B$ నిరూపణగా పేర్కొంది; చిత్రంలో $!A$ శాఖకు పొరపాటున $\delta_2$ గుర్తు పెట్టింది. $\delta_1$ను $!A$ శాఖకే, $\delta_2$ను $!A$ నుంచి $!B$ ఉపనిరూపణకే ఉంచాం.}

నిజానికి అలాంటి మార్గం ఎప్పుడూ ఉంటుంది. పైలాంటి "పక్కదారులు"—\tetoken{ఒక ప్రవేశ}{!!a{introduction}} నియమం వెంటనే \tetoken{ఒక తొలగింపు}{!!a{elimination}} నియమం రావడం—సహజ నిగమన \tetoken{నిరూపణల}{!!{proof}s} నుంచి ఎప్పుడూ తొలగించవచ్చు. దీన్ని \emph{సాధారణీకరణ} అంటారు; ఫలితం \emph{సాధారణ} \tetoken{నిరూపణ}{!!{proof}} (లేదా \emph{సాధారణ రూపంలో} గల \tetoken{నిరూపణ}{!!a{proof}}).

ఈ ఫలిత నిరూపణ సీక్వెంట్ తర్కంలోని కట్-తొలగింపు ప్రమేయంలానే కొంత సంక్లిష్టమైనది; అనేక సందర్భాలను పరిశీలించాలి. అంతఃప్రజ్ఞావాద \Log{N1i} కంటే సాంప్రదాయిక \Log{N1c}కు దీన్ని నిరూపించడం కష్టం. సీక్వెంట్ తర్కంలో \Cut{} నియమంతో అదనంగా ఏదీ \tetoken{నిరూపించలేమని}{!!{prove}} కట్ తొలగింపు చూపినట్లే, పక్కదారులను తప్పించుకున్నప్పటికీ వాటితో నిరూపించగల ఫలితాలన్నీ నిరూపించవచ్చని సాధారణీకరణ చూపుతుంది. ఇంకా సారూప్యాలు ఉన్నాయి: ఒకే ఫలితానికి పక్కదారులున్న అత్యల్ప \tetoken{నిరూపణ}{!!{proof}} కంటే సాధారణ \tetoken{నిరూపణ}{!!{proof}} ఎంతో పొడవుగా ఉండవచ్చు; సీక్వెంట్ తర్కంలో కట్‌లున్న \tetoken{నిరూపణలు}{!!{proof}s} కట్-రహిత అత్యల్ప \tetoken{నిరూపణ}{!!{proof}} కంటే ఎంతో చిన్నవిగా ఉండవచ్చునట్లే. తన \tetoken{నిరూపణలో}{!!{proof}} ఫలితపు ఉప-\tetoken{సూత్రాలనే}{!!{formula}s} ఉపయోగించి ఆ ఫలితాన్ని ఎప్పుడైనా \tetoken{నిరూపించవచ్చనే}{!!{prove}} ఉప-\tetoken{సూత్ర}{!!{formula}} లక్షణం, సీక్వెంట్ తర్కానికి కట్ తొలగింపు నుంచి వచ్చినట్లే సహజ నిగమనానికి సాధారణీకరణ నుంచి వస్తుంది.

\end{document}
